\documentclass{scrartcl}

\usepackage[hw]{darkwhite}

% cov
\DeclareMathOperator{\cov}{cov}

\title{18.676 Stochastic Calculus}
\subtitle{Problem Set 2}

\begin{document}
\section{Brownian Bridge}
\begin{prob*}
    We set $W_t = B_t - tB_1$ for every $t \in [0, 1]$. 
\end{prob*}
\begin{prob}
    Show that $(W_t)_{t \in [0, 1]}$ is a centered Gaussian process and give its covariance function.
\end{prob}
\begin{proof}
    We have that $W_t$ is a linear combination of Gaussian processes, and so is itself Gaussian. It is centered since $B_t$ is centered. The covariance function is given by
    \begin{align*}
        \cov(W_s, W_t) &= \cov(B_s - sB_1, B_t - tB_1) \\
        &= \cov(B_s, B_t) - s\cov(B_1, B_t) - t\cov(B_s, B_1) + st\cov(B_1, B_1) \\
        &= s\wedge t - st,
    \end{align*}
    where we have used the fact that $B_t$ is a Brownian motion, so that $\cov(B_s, B_t) = s\wedge t$.
\end{proof}
\begin{prob}
    Let $0 < t_1 < t_2 < \cdots < t_p < 1$. Show that the law of $(W_{t_1}, W_{t_2}, \ldots, W_{t_p})$ has density
    \[
        g(x_1, x_2, \ldots, x_p) = \sqrt{2\pi} p_{t_1}(x_1)p_{t_2 - t_1}(x_2 - x_1)\cdots p_{t_p - t_{p-1}}(x_p - x_{p-1})p_{1 - t_p}(-x_p),
    \]
    where $p_t(x) = \frac{1}{\sqrt{2\pi t}}e^{-x^2/2t}$, and explain why the law of $(W_{t_1}, W_{t_2}, \ldots, W_{t_p})$ can be interpreted as the conditional law of $(B_{t_1}, B_{t_2}, \ldots, B_{t_p})$ knowing that $B_1 = 0$.
\end{prob}

\begin{proof}
    Observe that $B_t - t B_1$ is independent of $B_1$ for all $t\in [0,1]$. Indeed,
    \begin{align*}
        \cov(B_t - tB_1, B_1) &= \cov(B_t, B_1) - t\cov(B_1, B_1) \\
        &= t - t \\
        &= 0.
    \end{align*}
    Thus, the density functions of $(W_{t_1},\dots, W_{t_p}, B_1)$ (denoted $f$) and $B_1$ (denoted $p_1$) will satisfy
    \[
        f(x_1,\dots, x_p, y) = g(x_1,\dots, x_p)p_1(y).
    \]
    \idea{
        Everything here is a GRV with support on its domain (I can draw any small ball in $\RR^{p+1}$ and there is some nonzero chance that $(W_{t_1},\dots, W_{t_p}, B_1)$ will lie in that ball, etc), so all of these densities exist.
    }
    However, it's pretty easy to compute $f$. Let $h$ be the density function of $(B_{t_1},\dots, B_{t_p}, B_1)$.
    \[
        f(x_1,\dots, x_p, y) = h(x_1 + t_1 y, \dots, x_p + t_p y, y).
    \]\pagebreak
    \idea{
        [The Jacobian is one]
        This linear transformation is the matrix
        \[
            \begin{pmatrix}
                1 & t_1 & 0 & \cdots & 0\\
                0 & 1 & t_2 & \cdots & 0\\
                \vdots & \vdots & \vdots & \ddots & \vdots\\
                0 & 0 & 0 & \cdots & 1
            \end{pmatrix},
        \]
        which has determinant one. So when we use change of variables, there's no Jacobian to worry about.
    }
    We know $h$ and $p_1$ though. Let $t_0 = 0$, $x_0 = 0$, $t_{p+1} = 1$, and $x_{p+1} = 0$.
    \begin{align*}
        h(x_1 + t_1 y,\dots, x_p + t_p y, y) &= \prod_{i=1}^{p+1} \frac{1}{\sqrt{2\pi (t_i - t_{i-1})}}e^{-((x_i + t_i y) - (x_{i-1} + t_{i-1}y))^2/2(t_i - t_{i-1})}
    \end{align*}
    Just stare at the exponents for a while. We have
    \begin{align*}
        \sum_{i=1}^{p+1} \frac{((x_i + t_i y) - (x_{i-1} + t_{i-1}y))^2}{2(t_i - t_{i-1})} &= \sum_{i=1}^{p+1} \frac{(x_i - x_{i-1})^2}{2(t_i - t_{i-1})}\\
        &+ 2y \sum_{i=1}^{p+1} \frac{(x_i - x_{i-1})(t_i - t_{i-1})}{2(t_i - t_{i-1})} + y^2\sum_{i=1}^{p+1} \frac{(t_i - t_{i-1})^2}{2(t_i - t_{i-1})} \\
        &= \sum_{i=1}^{p+1} \frac{(x_i - x_{i-1})^2}{2(t_i - t_{i-1})} + y^2.
    \end{align*}
    Thus,
    \[
        h(x_1 + t_1 y,\dots, x_p + t_p y, y) = \left(\prod_{i=1}^{p+1} \underbrace{\frac{1}{\sqrt{2\pi (t_i - t_{i-1})}}e^{-(x_i - x_{i-1})^2/2(t_i - t_{i-1})}}_{p_{t_i - t_{i-1}}(x_i - x_{i-1})}\right) p_1(y).
    \]
    and
    \[
        g(x_1,\dots, x_p) = \sqrt{2\pi} p_{t_1}(x_1)p_{t_2 - t_1}(x_2 - x_1)\cdots p_{t_p - t_{p-1}}(x_p - x_{p-1})p_{1 - t_p}(-x_p).
    \]
    as desired.
\end{proof}

% \idea{
%     Here is a massive computation, in case the first solution is not clear.
% }
% \begin{proof}[Solution 2 (Brute Force)]
%     We have that if $i < j$,
%     \begin{align*}
%         \cov(W_{t_i}, W_{t_j}) &= t_i - t_it_j \\
%         &= t_i(1 - t_j)
%     \end{align*}

%     Now, let $\Delta t_i = t_{i+1} - t_i$, and $\Delta W_i = W_{t_{i+1}} - W_{t_i}$ for $0\leq i \leq p$,
%     where we define $t_0 = 0$ and $t_{p+1} = 1$. Then, for $i < j$,
%     \begin{align*}
%         t_j\cov(\Delta W_i, \Delta W_j) &= \cov(W_{t_{i+1}} - W_{t_i}, W_{t_{j+1}} - W_{t_j}) \\
%         &= \cov(W_{t_{i+1}}, W_{t_{j+1}}) - \cov(W_{t_{i+1}}, W_{t_j}) - \cov(W_{t_i}, W_{t_{j+1}}) + \cov(W_{t_i}, W_{t_j}) \\
%         &= (t_{i+1} - t_{i+1}t_{j+1}) - (t_{i+1} - t_{i+1}t_j) - (t_i - t_it_{j+1}) + (t_i - t_it_j) \\
%         &= -t_{i+1}t_{j+1} + t_{i+1}t_j - t_it_{j+1} + t_it_j \\
%         &= -\Delta t_i \Delta t_j.
%     \end{align*}
%     % i.e. $\cov(\Delta W_i, \Delta W_j) = -1$.
%     Furthermore,
%     \begin{align*}
%         \cov(\Delta W_i, \Delta W_i) &= \cov(W_{t_{i+1}} - W_{t_i}, W_{t_{i+1}} - W_{t_i}) \\
%         &= \cov(W_{t_{i+1}}, W_{t_{i+1}}) - 2\cov(W_{t_{i+1}}, W_{t_i}) + \cov(W_{t_i}, W_{t_i}) \\
%         &= (t_{i+1} - t_{i+1}^2) - 2(t_{i+1} - t_{i+1}t_i) + (t_i - t_i^2) \\
%         &= (t_{i+1} - t_i) - (t_{i+1} - t_i)^2 \\
%         &= \Delta t_i - \Delta t_i^2.
%     \end{align*}
%     % i.e. $\cov(\Delta W_i, \Delta W_i) = \frac{1}{\Delta t_i} - 1$.
%     So, here's our covariance matrix between $\Delta W_0,\dots, \Delta W_{p-1}$:
%     \[
%         \Sigma = \begin{pmatrix}
%             \Delta t_0 - \Delta t_0^2 & -\Delta t_0 \Delta t_1 & \cdots & - \Delta t_0 \Delta t_{p-1} \\
%             -\Delta t_1 \Delta t_0 & \Delta t_1 - \Delta t_1^2 & \cdots & - \Delta t_1 \Delta t_{p-1} \\
%             \vdots & \vdots & \ddots & \vdots \\
%             -\Delta t_{p-1} \Delta t_0 & -\Delta t_{p-1} \Delta t_1 & \cdots & \Delta t_{p-1} - \Delta t_{p-1}^2
%         \end{pmatrix}.
%     \]
%     We're going to invert this guy! First off, observe that
%     \[
%         \Sigma \underbrace{\begin{pmatrix}
%             1 & 1 & \cdots & 1\\
%             1 & 1 & \cdots & 1\\
%             \vdots & \vdots & \ddots & \vdots\\
%             1 & 1 & \cdots & 1
%         \end{pmatrix}}_A = \Delta t_p\begin{pmatrix}
%             t_0 & t_0 & \cdots & t_0\\
%             t_1 & t_1 & \cdots & t_1\\
%             \vdots & \vdots & \ddots & \vdots\\
%             t_{p-1} & t_{p-1} & \cdots & t_{p-1}
%         \end{pmatrix}.
%     \]
%     Second, observe that
%     \[
%         \Sigma \underbrace{\begin{pmatrix}
%             \Delta t_0^{-1} & 0 & \cdots & 0\\
%             0 & \Delta t_1^{-1} & \cdots & 0\\
%             \vdots & \vdots & \ddots & \vdots\\
%             0 & 0 & \cdots & \Delta t_{p-1}^{-1}
%         \end{pmatrix}}_B = \begin{pmatrix}
%             1 - t_0 & - t_0 & \cdots & - t_0\\
%             - t_1 & 1 - t_1 & \cdots & - t_1\\
%             \vdots & \vdots & \ddots & \vdots\\
%             - t_{p-1} & - t_{p-1} & \cdots & 1 - t_{p-1}
%         \end{pmatrix}.
%     \]
%     Thus,
%     \[
%         \Sigma^{-1} = \frac{1}{\Delta t_p} A + B.  
%     \]
%     Great. Note that $(\Delta W_i)_{0\leq i\leq p-1}$ is a centered GRV because it is a linear combination of such. We know that the density of this variable is
%     \[
%         g(x_1, x_2, \ldots, x_p) = \frac{1}{Z}\exp\left(-\frac{1}{2}\sum_{i=1}^p\sum_{j=1}^p x_i(\Sigma^{-1})_{ij}x_j\right)
%     \]
%     where $Z$ is some normalizing constant. Given the expression for $\Sigma^{-1}$, this factors into
%     \[
%         \frac{1}{Z'}p_{\Delta t_0}(x_0) p_{\Delta t_1}(x_1) \cdots p_{\Delta t_{p-1}}(x_{p-1})\exp\left(-\frac{1}{2\Delta t_p}\left(\sum_{i=1}^p x_i\right)^2\right).
%     \]
%     where $Z'$ is some other constant.
%     (This is a valid probability distribution, thus the constant $Z$ has been absorbed). Finally, changing variables back to $B_1,\dots, B_p$ via a linear transformation with constant Jacobian, we get the desired density.
    

    
    % where we have set $t_0 = 0$. This is the covariance matrix of a multivariate Gaussian distribution with mean zero, and so the law of $(W_{t_1}, W_{t_2}, \ldots, W_{t_p})$ is Gaussian. The density is given by
    % \begin{align*}
    %     g(x_1, x_2, \ldots, x_p) &= \frac{1}{\sqrt{\det(2\pi\Sigma)}}\exp\left(-\frac{1}{2}(x - \mu)^T\Sigma^{-1}(x - \mu)\right) \\
    %     &= \frac{1}{\sqrt{2\pi}^p\sqrt{\det(\Sigma)}}\exp\left(-\frac{1}{2}\sum_{i=1}^p\sum_{j=1}^p(x_i - \mu_i)(\Sigma^{-1})_{ij}(x_j - \mu_j)\right),
    % \end{align*}
    % where $\mu = 0$ and $\Sigma_{ij} = \cov(W_{t_i}, W_{t_j})$. We have that
    % \begin{align*}
    %     (\Sigma^{-1})_{ij} &= \frac{1}{\det(\Sigma)}\left(\cov(W_{t_i}, W_{t_j})\right)^{ij} \\
    %     &= \frac{1}{\det(\Sigma)}\left((t_i - t_{i-1})\wedge(t_j - t_{j-1}) - (t_i - t_{i-1})(t_j - t_{j-1})\right),
    % \end{align*}
    % and so
    % The law of $(W_{t_1}, W_{t_2}, \ldots, W_{t_p})$ can be interpreted as the conditional law of $(B_{t_1}, B_{t_2}, \ldots, B_{t_p})$ knowing that $B_1 = 0$ since $W_t = B_t - tB_1$.
% \end{proof}
\begin{prob}
    Verify that the two processes $(W_t)_{t \in [0, 1]}$ and $(W_{1-t})_{t \in [0, 1]}$ have the same distribution (similarly as in the definition of Wiener measure, this law is a probability measure on the space of all continuous functions from $[0, 1]$ into $\RR$).
\end{prob}
\begin{proof}
    The density we just showed is manifestly time-reversal-invariant. Of course, prescribing this density is equivalent to specifying the law on all finite-dimensional marginals, i.e. the probability measure over any cylinder set consisting of continuous functions which pass through Borel sets $A_1,\dots A_n$ at times $t_1 < \dots < t_n$. These cylinder sets form a generating $\pi$-system for the product $\sigma$-algebra, and thus the law of $(W_t)_{t\in [0,1]}$ is time-reversal-invariant.
\end{proof}
\section{Time Reversal for Brownian Motion}
\begin{prob*}
    We set $B_0^t = B_1B_1^t$ for every $t \in [0, 1]$. Show that the two processes $(B_t)_{t \in [0, 1]}$ and $(B_0^t)_{t \in [0, 1]}$ have the same law (as in the definition of Wiener measure, this law is a probability measure on the space of all continuous functions from $[0, 1]$ into $\RR$).
\end{prob*}
\idea{
    [Hint]
    Note that the choice of end time $t = 1$ is not special: it will follow by essentially the same argument that for any fixed $t \geq 0$, the process $(B_t - B_{t-s}: 0 \leq s \leq t)$ is a standard Brownian motion on the time interval $[0, t]$.
}
\section{Reflected Brownian Motion}
\begin{prob*}
    Let $B_t$ be a standard Brownian motion, $S_t = \sup\{B_s: 0 \leq s \leq t\}$ its running maximum, and $Y_t = S_t - B_t \geq 0$. In this exercise you will show that the processes $(Y_t)_{t \geq 0}$ and $(\abs{B_t})_{t \geq 0}$ are equal in law. Note: be careful that $(S_t)_{t \geq 0}$ is in fact not a Markov process, although it will follow from this exercise that $(Y_t)_{t \geq 0}$ is.
\end{prob*}
\begin{prob}
    First argue that to prove the claim, it suffices to show that for any fixed $s, t \geq 0$, conditional on $\FFF_s$ the random variable $Y_{s+t}$ is equidistributed as $\abs{Y_s + \bar{B}_t}$ where $\bar{B}$ denotes a Brownian motion independent of $\FFF_s$.
\end{prob}

What it means for $(Y_t)_{t\geq 0}$ and $(\abs{B_t})_{t\geq 0}$ to have the same law is that for all $t_1, t_2, \ldots, t_n \geq 0$, the joint distribution of $(Y_{t_1}, Y_{t_2}, \ldots, Y_{t_n})$ is the same as the joint distribution of $(\abs{B_{t_1}}, \abs{B_{t_2}}, \ldots, \abs{B_{t_n}})$. The key is to break this up into conditional distributions, and then use the fact that the law of $(B_t)_{t\geq 0}$ is invariant under time reversal.

Start with the expression
\[
    \PP\left(Y_{s+t} \in A \mid \FFF_s\right) = \PP\left(\abs{Y_s + \bar{B}_t} \in A \mid \FFF_s\right), \quad a \geq 0,\qquad A \in \BBB(\RR).
\]
We will now show by induction that
\[
    \PP(\bigcap_i (Y_{t_i}\in A_i)) = \PP(\bigcap_i (\abs{B_{t_i}}\in A_i))
\]
for all $t_1, t_2, \ldots, t_n \geq 0$ and $A_1, A_2, \ldots, A_n \in \BBB(\RR)$. The base case is $n = 0$ which is trivial.

Then, as the inductive step,
\begin{align*}
    \PP\left(\bigcap_i (Y_{t_i}\in A_i)\right) &= \PP\left(\left(\abs{Y_{t_{i - 1} +\bar{B}_{t_i - t_{i-1}}}}\in A_i\right)\cap \bigcap_{j < i} (Y_{t_j}\in A_j)\right)\\
    &= \PP\left(\left(\abs{B_{t_{i - 1} +\bar{B}_{t_i - t_{i-1}}}}\in A_i\right)\cap \bigcap_{j < i} (B_{t_j}\in A_j)\right)\\
    &= \PP\left(\bigcap_i (\abs{B_{t_i}}\in A_i)\right).
\end{align*}
where the last line is because conditioned on $\{B_{t_j}: j < i\}$, $B_{t_{i - 1}} + \bar{B}_{t_i - t_{i-1}}$ has the same law as $B_{t_{i - 1} + B_{t_i - t_{i-1}}} = B_{t_i}$.

The inductive hypothesis is precisely the claim that $(Y_t)_{t\geq 0}$ and $(\abs{B_t})_{t\geq 0}$ have the same law (we have verified it over all finite-dimensional cylinder sets, which form a $\pi$-system generating the product $\sigma$-algebra). Thus, we have shown that the two processes have the same law.

\begin{prob}
    For $u \geq 0$ let $\hat{B}_u \equiv B_{s+u} - B_s$ and $\hat{S}_u = \sup\{\hat{B}_w: 0 \leq w \leq u\}$. Show that $Y_{s+t} = \max\{Y_s, \hat{S}_t\} - \hat{B}_t$.
\end{prob}
\begin{prob}
    Argue that for any fixed $y \geq 0$, $\max\{y, \hat{S}_t\} - \hat{B}_t$ has the same law as $\abs{y + \hat{B}_t}$.
\end{prob}
\idea{
    [Hint]
    A hint is to first decompose
    \[
        \PP\left(\max\{y, \hat{S}_t\} - \hat{B}_t > a\right) = \PP\left(y - \hat{B}_t > a\right) + \PP\left(y - \hat{B}_t \leq a, \hat{S}_t - \hat{B}_t > a\right),
    \]
    and use time reversal to simplify the last term.
}
\section{Arcsine Law}
\begin{prob*}
    Set $T = \inf\{t \geq 0: B_t = S_1\}$.
\end{prob*}
\begin{prob}
    Show that $T < 1$ a.s. (one may use the result of the previous exercise) and then that $T$ is not a stopping time.
\end{prob}

We know the Brownian motion is continuous almost surely. Continuous functions on compact domains attain their maximum, thus there exists a $t\in [0,1]$ such that $B_t = S_1$ almost surely. It remains to show that $B_1 = S_1$ almost never. To do this, we use time reversal; let $B'_t = B_{1-t} - B_1$, and note that $B'_t$ is a Brownian motion; then $S'_1 = \inf\{B'_t, 0 \leq t \leq 1\} = S_1 - B_1$. However, as a Brownian motion, we know that almost surely for every $\eps > 0$,
\[
    \sup_{0 \leq t \leq \eps} B'_t > 0
\]
thus $S'_1 > 0$ almost surely, i.e. $S_1 > B_1$ almost surely. So we conclude.


\begin{prob}
    Verify that the three variables $S_t$, $S_t - B_t$, and $\abs{B_t}$ have the same law.
\end{prob}
\begin{prob}
    Show that $T$ is distributed according to the so-called arcsine law, whose density is
    \[
        g(t) = \frac{1}{\pi\sqrt{t(1-t)}} \id_{(0,1)}(t).
    \]
\end{prob}
\begin{prob}
    Show that the results of questions 1. and 3. remain valid if $T$ is replaced by $L = \sup\{t \leq 1: B_t = 0\}$.
\end{prob}
\idea{
    [Hint]
    You may use the result of Problem 3 in your solution; it is especially useful for the last part of this exercise.
}
\section{Law of the Iterated Logarithm}
\begin{prob*}
    The goal of the exercise is to prove that
    \[
        \limsup_{t \to \infty} \frac{B_t}{\sqrt{2t\log\log t}} = 1 \quad \text{a.s.}
    \]
    We set $h(t) = \sqrt{2t\log\log t}$.
\end{prob*}
\idea{
    [Hint]
    In your solution, please give a proof of the following standard Gaussian tail bound: if $Z \sim \mathcal{N}(0, 1)$, then for all $x > 0$ we have
    \[
        \left(1 - \frac{1}{x^2}\right)\exp\left(-\frac{x^2}{2}\right) \leq \PP(Z \geq x) \leq \exp\left(-\frac{x^2}{2}\right)\frac{1}{x\sqrt{2\pi}}.
    \]
}
\begin{prob}
    Show that, for every $t > 0$, $\PP(S_t > u\sqrt{t}) \sim \frac{2}{u\sqrt{2\pi}} e^{-u^2/2}$ when $u \to +\infty$.
\end{prob}
\begin{prob}
    Let $r$ and $c$ be two real numbers such that $1 < r < c^2$. From the behavior of the probabilities $\PP(S_{r^n} > c h(r^{n-1}))$ when $n \to 1$, infer that, a.s.,
    \[
        \limsup_{t \to \infty} \frac{B_t}{h(t)} \geq 1.
    \]
\end{prob}
\begin{prob}
    Show that a.s. there are infinitely many values of $n$ such that
    \[
        B_{r^n} - B_{r^{n-1}} \geq \sqrt{\frac{r-1}{r}} h(r^{n-1}).
    \]
    Conclude that the statement given at the beginning of the exercise holds.
\end{prob}
\begin{prob}
    What is the value of $\liminf_{t \to \infty} \frac{B_t}{h(t)}$?
\end{prob}


% Exercise 2.27 (Brownian bridge)
% We set Wt D Bt  tB1 for every t 2 Œ0; 1.
% 1. Show that .Wt/t2Œ0;1 is a centered Gaussian process and give its covariance
% function.
% 2. Let 0 < t1 < t2 <  < tp < 1. Show that the law of .Wt1 ; Wt2 ;:::; Wtp / has
% density
% g.x1;:::; xp/ D p2 pt1 .x1/pt2t1 .x2  x1/ ptptp1 .xp  xp1/p1tp .xp/;
% where pt.x/ D p
% 1
% 2t
% exp.x2=2t/. Explain why the law of.Wt1 ; Wt2 ;:::; Wtp / can
% be interpreted as the conditional law of .Bt1 ; Bt2 ;:::; Btp / knowing that B1 D 0.
% Exercises 39
% 3. Verify that the two processes .Wt/t2Œ0;1 and .W1t/t2Œ0;1 have the same distribution (similarly as in the definition of Wiener measure, this law is a probability
% measure on the space of all continuous functions from Œ0; 1 into R).
% 1. LG Exercise 2.27 (Brownian bridge).
% 2. LG Exercise 2.31 (time reversal for Brownian motion). Note that the choice of end time 𝑡 = 1 is not special:
% it will follow by essentially the same argument that for any fixed 𝑡 ≥ 0, the process (𝐵𝑡 − 𝐵𝑡−𝑠
% : 0 ≤ 𝑠 ≤ 𝑡) is a
% standard Brownian motion on the time interval [0, 𝑡].
% 3. Reflected Brownian motion. Let 𝐵𝑡 be a standard Brownian motion, 𝑆𝑡 = sup{𝐵𝑠
% : 0 ≤ 𝑠 ≤ 𝑡} its running
% maximum, and 𝑌𝑡 = 𝑆𝑡 − 𝐵𝑡 ≥ 0. In this exercise you will show that the processes (𝑌𝑡)𝑡≥0 and (|𝐵𝑡
% |)𝑡≥0 are equal
% in law. Note: be careful that (𝑆𝑡)𝑡≥0 is in fact not a Markov process, although it will follow from this exercise that
% (𝑌𝑡)𝑡≥0 is.
% (a) First argue that to prove the claim, it suffices to show that for any fixed 𝑠, 𝑡 ≥ 0, conditional on ℱ𝑠 the
% random variable 𝑌𝑠+𝑡
% is equidistributed as |𝑌𝑠 + 𝐵¯
% 𝑡
% | where 𝐵¯ denotes a Brownian motion independent of ℱ𝑠
% .
% (b) For 𝑢 ≥ 0 let 𝐵ˆ
% 𝑢 ≡ 𝐵𝑠+𝑢 − 𝐵𝑠 and 𝑆ˆ
% 𝑢 ≡ sup{𝐵ˆ𝑤 : 0 ≤ 𝑤 ≤ 𝑢}. Show that 𝑌𝑠+𝑡 = max{𝑌𝑠
% , 𝑆ˆ
% 𝑡 } − 𝐵ˆ
% 𝑡
% .
% (c) Argue that for any fixed 𝑦 ≥ 0, max{𝑦, 𝑆ˆ
% 𝑡 } − 𝐵ˆ
% 𝑡 has the same law as |𝑦 + 𝐵ˆ
% 𝑡
% |. A hint is to first decompose
% ℙ
% 
% max{𝑦, 𝑆ˆ
% 𝑡 } − 𝐵ˆ
% 𝑡 > 𝑎
% 
% = ℙ
% 𝑦 − 𝐵ˆ
% 𝑡 > 𝑎
% 
% + ℙ
% 
% 𝑦 − 𝐵ˆ
% 𝑡 ≤ 𝑎, 𝑆ˆ
% 𝑡 − 𝐵ˆ
% 𝑡 > 𝑎
% 
% , (1)
% and use time reversal (Problem 2) to simplify the last term.
% 4. LG Exercise 2.32 (arcsine law). Hint: you may use the result of Problem 3 in your solution; it is especially useful
% for the last part of this exercise.
% 5. LG Exercise 2.33 (law of the iterated logarithm). In your solution, please give a proof of the following
% standard gaussian tail bound: if 𝑍 ∼ 𝒩(0, 1), then for all 𝑥 > 0 we have
% 
% 1 −
% 1
% 𝑥
% 2
% 
% exp(−𝑥
% 2
% /2)
% 𝑥
% √
% 2𝜋
% ≤ ℙ(𝑍 ≥ 𝑥) ≤ exp(−𝑥
% 2
% /2)
% 𝑥
% √
% 2𝜋
% . (2)
% Exercise 2.31 (Time reversal) We set B0
% t D B1B1t for every t 2 Œ0; 1. Show that
% the two processes .Bt/t2Œ0;1 and .B0
% t
% /t2Œ0;1 have the same law (as in the definition of
% Wiener measure, this law is a probability measure on the space of all continuous
% functions from Œ0; 1 into R).
% Exercise 2.32 (Arcsine law)
% Set T WD infft  0 W Bt D S1g:
% 1. Show that T < 1 a.s. (one may use the result of the previous exercise) and then
% that T is not a stopping time.
% 2. Verify that the three variables St, St  Bt and jBtj have the same law.
% 3. Show that T is distributed according to the so-called arcsine law, whose density
% is
% g.t/ D 1

% pt.1  t/
% 1.0;1/.t/:
% 4. Show that the results of questions 1. and 3. remain valid if T is replaced by
% L WD supft  1 W Bt D 0g:
% Exercise 2.33 (Law of the iterated logarithm)
% The goal of the exercise is to prove that
% lim sup
% t!1
% Bt p2t log log t
% D 1 a.s.
% 40 2 Brownian Motion
% We set h.t/ D p2t log log t:
% 1. Show that, for every t > 0, P.St > u
% pt/  2
% u
% p2
% eu2=2, when u ! C1.
% 2. Let r and c be two real numbers such that 1 < r < c2: From the behavior of the
% probabilities P.Srn > c h.rn1// when n ! 1, infer that, a.s.,
% lim sup
% t!1
% Bt p2t log log t
%  1:
% 3. Show that a.s. there are infinitely many values of n such that
% Brn  Brn1
% rr  1
% r h.r
% n
% /:
% Conclude that the statement given at the beginning of the exercise holds.
% 4. What is the value of lim inf t!1
% Bt p2t log log t
% ?

\end{document}
